\section{\textbf{Conclusions}}
This paper has presented consciousness not as a single property awaiting spontaneous emergence, but as a set of distinguishable capacities that can be constructed, examined, and tested. By separating phenomenal, observable, and non-emergent consciousness, the framework preserves the unresolved problem of subjective experience while enabling rigorous investigation of integrated world modeling, self-representation, memory, metacognition, temporal continuity, and adaptive agency. Language serves within this account as a unifying semantic interface that makes heterogeneous information mutually accessible without replacing modality-specific representations.
The proposed testing framework further distinguishes internal representation, global access, and explicit report through convergent evidence, causal intervention, dissociation, anti-imitation controls, and longitudinal evaluation. Success under these tests would support claims about observable functional consciousness, but would not prove phenomenal experience. This approach therefore replaces reliance on scale or an unknown threshold with an experimentally grounded program for designing and evaluating increasingly unified artificial minds.